\section{INTRODUCCIÓN}

Los autómatas celulares son una 4-tupla que se define como:

\begin{itemize}
	\item Un espacio $D$-dimensional en donde se encuentran agrupadas un conjunto de células
	\item Un conjunto de estados, cada célula puede adquirir uno de esos estados
	\item Un conjunto de vectores de desplazamiento que definen la vecindad de cada célula
	\item Una función de transición que define el comportamiento general.
\end{itemize}

Todas las células utilizan la misma función de transición de forma simultanea, esta función depende de los estados de las vecinas, es decir, de aquellas células que se encuentran en su vector de desplazamiento.

Estos sistemas tienen la capacidad de simular sistemas complejos, el comportamiento usual para hacerlo es definiendo todos los elementos de la tupla del autómata y después aplicarlos a una condición inicial para iniciar su evolución. Sin embargo, este proceso puede llegar a ser lento debido a la complejidad y la fidelidad del evento. Es por esto que se propone un método inverso, dado un conjunto de datos, definir el función de transición para poder extrapolar el comportamiento a nuevas condiciones iniciales. 

Para esto, se utiliza una red neuronal que se entrena con la evolución de un sistema complejo. La red entrenada se sustituye por la función de transición del autómata celular dando vida a los autómatas celulares neuronales.

La definición de estos autoamtas celulares, a diferencia de los estudiados por Mordvinstev et. al. en \cite{}, se enfocan a espacios bidimensionales, discretos y co nestados escalares (no vectores con estados ocultos), de forma determinista. Esto para probar que es posible con una definición mas pequeña obtener comportamientos similares. 

Se realizarón 6 experimentos, 3 sobre simulaciones meramente matematicos (el juego de la vida, el \textit{sandpile} y la hormiga de Langton) y 3 sobre sistemas complejos reales: propagación de fuego (FLAME3), movimiento de tormentas (SEVIR) y el crecimiento del moho \textit{Physarum polycephalum} en una caja de Petri.

Las simulaciones en los sistemas bien conocidos permitieron definir la composición y los metaparametros de la red neuronal. Se pudieron replicar todas las configuraciones tradicionales de cada experimento, se encontro una posible configuración estable del juego de la vida, se experimentaron limites de multiples generadore sen el \textit{sandpile} y se experiemnto con multples hormigas dentro del mismo espacio con multiples estados.

La alicacion en los sistemas reales, de los cuales solo FLAME3 y  \textit{Physarum polycephalum} pudieron replicarse, mientras que se pudo simular un comportamiento de SEVIR pero se encontro limitado por ser un sistema abierto. 